\topic{Paradoks Pemungutan Suara}
\index{paradoks pemungutan suara|(}
%\emph{(Optional material from this chapter is discussed here,
%but is not required to complete the exercises.)}

Bayangkan bahwa sebuah kelas Ilmu Politik yang mempelajari proses
pemilihan presiden Amerika mengadakan simulasi pemilihan.
Sebanyak $29$ anggota kelas
mengurutkan calon Partai Demokrat,
Partai Republik, dan Partai Ketiga,
dari yang paling disukai hingga yang paling tidak disukai
($>$ berarti `lebih disukai daripada').
\begin{center}
  \renewcommand{\arraystretch}{1.1}
  \begin{tabular}{cc}
    \multicolumn{1}{c}{\begin{tabular}[b]{@{}c@{}}
       \textit{urutan preferensi}
    \end{tabular}}
    &\multicolumn{1}{l}{{\renewcommand{\arraystretch}{1.0}
      \begin{tabular}[b]{@{}l@{}}
         \textit{jumlah pemilih dengan}    \\[-0.5ex]
         \textit{preferensi tersebut}
      \end{tabular}}}                \\  
    \hline
    \begin{tabular}{@{}r@{}}
      $\text{Demokrat} > \text{Republik} > \text{Ketiga}$      \\
      $\text{Demokrat} > \text{Ketiga} > \text{Republik}$      \\
      $\text{Republik} > \text{Demokrat} > \text{Ketiga}$      \\
      $\text{Republik} > \text{Ketiga} > \text{Demokrat}$      \\
      $\text{Ketiga} > \text{Demokrat} > \text{Republik}$      \\
      $\text{Ketiga} > \text{Republik} > \text{Demokrat}$   
    \end{tabular}
    &\begin{tabular}{@{}r@{}}                                                
      $5$     \\
      $4$     \\
      $2$     \\
      $8$     \\
      $8$     \\
      $2$    
    \end{tabular}    
  \end{tabular}
    \renewcommand{\arraystretch}{1}
\end{center}
Bagaimanakah preferensi kelompok itu secara keseluruhan?

Secara keseluruhan,
kelompok tersebut lebih menyukai calon Demokrat daripada calon Republik dengan selisih lima suara;
tujuh belas pemilih menempatkan calon Demokrat di atas calon Republik,
sedangkan dua belas pemilih memilih urutan sebaliknya.
Kelompok tersebut juga
lebih menyukai calon Republik daripada calon Partai Ketiga,
lima belas berbanding empat belas.
Namun, anehnya, kelompok itu juga
lebih menyukai calon Partai Ketiga daripada calon Demokrat, delapan belas berbanding sebelas.
\begin{center}
  \includegraphics{vs/mp/voting.1}
\end{center}
Ini adalah sebuah \definend{paradoks pemungutan suara}\index{paradoks pemungutan suara!definisi},
lebih khusus lagi, sebuah
\definend{siklus mayoritas}.\index{paradoks pemungutan suara!siklus mayoritas}

Matematikawan mempelajari paradoks pemungutan suara antara lain karena
implikasinya bagi politik praktis.
Sebagai contoh, pengajar dapat memanipulasi kelas ini
agar memilih calon Demokrat sebagai pemenang keseluruhan
dengan mula-mula meminta pemungutan suara antara calon Republik dan
calon Partai Ketiga,
lalu meminta pemungutan suara antara pemenang kontes itu, yaitu calon
Republik, dan calon Demokrat.
Dengan manipulasi serupa,
pengajar dapat membuat salah satu dari dua calon lainnya menjadi pemenang.
(Kita akan membatasi pembahasan pada pemilihan dengan tiga calon, tetapi hal yang sama
terjadi dalam pemilihan yang lebih besar.)

Matematikawan juga mempelajari paradoks pemungutan suara semata-mata karena
paradoks itu menarik.
Salah satu aspek yang menarik ialah bahwa
siklus mayoritas kelompok secara keseluruhan terjadi meskipun
daftar preferensi setiap pemilih bersifat
\definend{rasional}\index{paradoks pemungutan suara!urutan preferensi rasional}, yakni
berupa urutan lurus.
Dengan kata lain, siklus mayoritas tampaknya muncul dalam keseluruhan
tanpa hadir dalam komponen-komponen keseluruhan tersebut, yaitu daftar preferensi.
Namun, kita dapat menggunakan aljabar linear untuk menunjukkan bahwa
kecenderungan menuju preferensi siklik sebenarnya terdapat
dalam daftar setiap pemilih dan
bahwa kecenderungan itu tampak ketika penjumlahannya lebih besar
daripada pembatalannya.

Untuk itu,
dengan menyingkat pilihan sebagai $D$, $R$, dan $T$,
kita dapat menggambarkan bagaimana
seorang pemilih dengan urutan preferensi $D>R>T$ berkontribusi terhadap siklus di atas.
\begin{center}
  \includegraphics{vs/mp/voting.2}
\end{center}
(Tanda negatif muncul karena panah itu menggambarkan $T$ sebagai
lebih disukai daripada $D$, sedangkan pemilih ini menyukai urutan sebaliknya.)
Gambaran untuk daftar preferensi lainnya terdapat pada tabel di
halaman~\pageref{table:Voting}.

Sekarang, untuk melaksanakan pemilihan, kita menggabungkan gambaran-gambaran ini secara linear;
sebagai contoh, simulasi pemilihan kelas Ilmu Politik
\begin{equation*}
  5\cdot \votinggraphic{43}
  +4\cdot\votinggraphic{44}
  +\dots+2\cdot\votinggraphic{45}
\tag*{}\end{equation*}
menghasilkan preferensi kelompok melingkar yang ditunjukkan sebelumnya.

Tentu saja, membentuk kombinasi linear merupakan aljabar linear.
Notasi siklus grafis memberi gambaran yang baik, tetapi kurang praktis, sehingga kita
menggunakan vektor kolom dengan memulai dari $D$ dan
mengambil bilangan-bilangan pada siklus dengan urutan berlawanan arah jarum jam.
Dengan demikian, simulasi pemilihan dan satu suara $D>R>T$ kita nyatakan sebagai berikut.
\begin{equation*}
  \votepreflist{7}{1}{5}
  \quad\text{dan}\quad
  \votepreflist{-1}{1}{1}
\tag*{}\end{equation*}

Kita akan
menguraikan vektor suara menjadi dua bagian, yang satu
siklik dan yang lain asiklik.
Untuk bagian pertama,
kita mengatakan bahwa sebuah vektor \definend{siklik murni} jika vektor itu berada dalam
subruang $\Re^3$ berikut.
\begin{equation*}
  C
  =\set{\votepreflist{k}{k}{k} \suchthat k \in\Re}
  =\set{k\cdot\votepreflist{1}{1}{1} \suchthat k \in\Re}
\end{equation*}
Untuk bagian kedua, perhatikan himpunan
vektor yang
tegak lurus terhadap semua vektor dalam $C$.
\nearbyexercise{exer:PerpIsSubSp} menunjukkan bahwa himpunan ini merupakan subruang
\begin{align*}
  C^\perp &=\set{\votepreflist{c_1}{c_2}{c_3} \suchthat 
                 \colvec{c_1 \\ c_2 \\ c_3}\dotprod\colvec{k \\ k \\ k} =0
                 \text{\ untuk setiap $k\in\Re$}}             \\
          &=\set{\votepreflist{c_1}{c_2}{c_3} \suchthat 
                 c_1+c_2+c_3=0}                                         
          =\set{c_2\votepreflist{-1}{1}{0}
                +c_3\votepreflist{-1}{0}{1} \suchthat  c_2,c_3\in\Re}
\end{align*}
(dibaca ``$C$~tegak lurus'').
Dengan demikian, kita memperoleh basis berikut untuk $\Re^3$.
\begin{equation*}
 \sequence{\votepreflist{1}{1}{1},
             \votepreflist{-1}{1}{0},
             \votepreflist{-1}{0}{1}}
\tag*{}\end{equation*}
Kita dapat menyatakan suara terhadap basis ini,
dan dengan demikian menguraikannya menjadi bagian siklik dan bagian asiklik.
\textit{(Catatan bagi pembaca yang telah mempelajari bagian pilihan
dalam bab ini:~artinya, ruang tersebut merupakan jumlah langsung dari $C$ dan~$C^\perp$.)}

Sebagai contoh, perhatikan pemilih $D>R>T$ yang dibahas di atas.
Kita menyatakan pemilih tersebut terhadap basis ini melalui
\begin{equation*}
  \begin{linsys}{3}
    c_1  &-  &c_2  &-  &c_3  &=  &-1 \\
    c_1  &+  &c_2  &   &     &=  &1  \\
    c_1  &   &     &+  &c_3  &=  &1  
  \end{linsys}
  \grstep[-\rho_1+\rho_3]{-\rho_1+\rho_2}
  \repeatedgrstep{(-1/2)\rho_2+\rho_3}
  \begin{linsys}{3}
    c_1  &-  &c_2  &-  &c_3      &=  &-1 \\
         &   &2c_2 &+  &c_3      &=  &2  \\
         &   &     &   &(3/2)c_3 &=  &1  
  \end{linsys}
\end{equation*}
dengan menggunakan koordinat $c_1=1/3$, $c_2=2/3$, dan $c_3=2/3$.
Selanjutnya,
\begin{equation*}
  \votepreflist{-1}{1}{1}
  =\frac{1}{3}\cdot\votepreflist{1}{1}{1}
    +\frac{2}{3}\cdot\votepreflist{-1}{1}{0}
    +\frac{2}{3}\cdot\votepreflist{-1}{0}{1}
  =\votepreflist{1/3}{1/3}{1/3}
   +\votepreflist{-4/3}{2/3}{2/3}
\tag*{}\end{equation*}
memberikan penguraian yang diinginkan menjadi
bagian siklik dan bagian asiklik.
\begin{equation*}
  \votinggraphic{3}
=
\votinggraphic{4}
\>+\>
\votinggraphic{5}
\end{equation*}
Dengan demikian, kita melihat bahwa daftar preferensi rasional
pemilih $D>R>T$ ini memang memiliki bagian siklik.

Pemilih $T>R>D$ berlawanan dengan pemilih yang baru saja dibahas
karena simbol `$>$' dibalik.
Penguraian pemilih ini
\begin{equation*}
\votinggraphic{6}
=
\votinggraphic{7}
\>+\>
\votinggraphic{8}
\end{equation*}
menunjukkan bahwa preferensi yang berlawanan ini mempunyai penguraian yang juga berlawanan.
Kita mengatakan bahwa pemilih pertama mempunyai
\definend{putaran} positif\index{putaran}\index{paradoks pemungutan suara!putaran}
karena bagian siklusnya searah dengan
arah panah yang telah kita pilih,
sedangkan putaran pemilih kedua bernilai negatif.

Fakta bahwa kedua pemilih yang berlawanan ini saling membatalkan tercermin dalam
kenyataan bahwa jumlah vektor suaranya adalah nol.
Hal ini menyarankan cara lain untuk menghitung hasil pemilihan.
Kita dapat terlebih dahulu membatalkan sebanyak mungkin pasangan daftar preferensi yang berlawanan,
kemudian menentukan hasilnya dengan menjumlahkan daftar yang tersisa.

Tabel di bawah memuat tiga pasangan
daftar preferensi yang berlawanan.
Sebagai contoh, baris teratas memuat para pemilih yang dibahas di atas.
\begin{center}
  \def\betweenrowvspace{5.5ex}
  \makebox[\hsize][l]{%  box runs into margin; suppress overfull warn
  \begin{tabular}{@{}c@{\hspace{0em}}c@{}} 
    % \hline
    \multicolumn{1}{c}{\textit{putaran positif}}
    &\multicolumn{1}{c}{\textit{putaran negatif}}  
     \\ \hline
    {\renewcommand{\arraystretch}{1.4}%%
     \begin{tabular}{@{}c}
      $\text{Demokrat}>\text{Republik}>\text{Ketiga}$                     \\
      $\votinggraphic{9} % \votecycle{1}{1}{-1}{3}
      =\votinggraphic{10} %\votecycle{1/3}{1/3}{1/3}{3}
      \>+\>
      \votinggraphic{11}$  %\votecycle{2/3}{2/3}{-4/3}{3}\;{}
    \end{tabular}}
    &{\renewcommand{\arraystretch}{1.4}%%
     \begin{tabular}{c@{}}
      $\text{Ketiga}>\text{Republik}>\text{Demokrat}$                     \\
      $\votinggraphic{12}  %\votecycle{-1}{-1}{1}{3}
      =\votinggraphic{13}  %\votecycle{-1/3}{-1/3}{-1/3}{3}
      \>+\>
      \votinggraphic{14}$  %\votecycle{-2/3}{-2/3}{4/3}{3}\;{}
    \end{tabular}}                  \\[\betweenrowvspace] % \hline
    {\renewcommand{\arraystretch}{1.4}%%
    \begin{tabular}{@{}c}
      $\text{Republik}>\text{Ketiga}>\text{Demokrat}$                     \\
      $\votinggraphic{15}  %\votecycle{-1}{1}{1}{3}
      =\votinggraphic{16} %\votecycle{1/3}{1/3}{1/3}{3}
      \>+\>
      \votinggraphic{17}$  %\votecycle{-4/3}{2/3}{2/3}{3}\;{}
    \end{tabular}}
    &{\renewcommand{\arraystretch}{1.4}%%
     \begin{tabular}{c@{}}
      $\text{Demokrat}>\text{Ketiga}>\text{Republik}$                     \\
      $\votinggraphic{18}  %\votecycle{1}{-1}{-1}{3}
      =\votinggraphic{19}  %\votecycle{-1/3}{-1/3}{-1/3}{3}
      \>+\>
      \votinggraphic{20}$  %\votecycle{4/3}{-2/3}{-2/3}{3}\;{}
    \end{tabular}}                  \\[\betweenrowvspace] % \hline
    {\renewcommand{\arraystretch}{1.4}%%
    \begin{tabular}{@{}c}
      $\text{Ketiga}>\text{Demokrat}>\text{Republik}$                     \\
      $\votinggraphic{21}  %\votecycle{1}{-1}{1}{3}
      =\votinggraphic{22}  %\votecycle{1/3}{1/3}{1/3}{3}
      \>+\>
      \votinggraphic{23}$  %\votecycle{2/3}{-4/3}{2/3}{3}\;{}
    \end{tabular}}
    &{\renewcommand{\arraystretch}{1.4}%%
     \begin{tabular}{c@{}}
      $\text{Republik}>\text{Demokrat}>\text{Ketiga}$                     \\
      $\votinggraphic{24}  %\votecycle{-1}{1}{-1}{3}
      =\votinggraphic{25}  %\votecycle{-1/3}{-1/3}{-1/3}{3}
      \>+\>
      \votinggraphic{26}$  %\votecycle{-2/3}{4/3}{-2/3}{3}\;{}
    \end{tabular}}                  % \\ \hline
  \end{tabular}
  }
  \label{table:Voting}
\end{center}
Jika pemilihan dilaksanakan seperti yang baru saja dijelaskan, maka setelah kita
membatalkan sebanyak mungkin pasangan pemilih yang berlawanan, akan tersisa tiga
himpunan daftar preferensi:~satu himpunan dari baris pertama, satu dari baris kedua,
dan satu dari baris ketiga.
Sebagai penutup, kita akan membuktikan bahwa
paradoks pemungutan suara hanya dapat terjadi jika
putaran ketiga himpunan itu searah.
Artinya, agar paradoks pemungutan suara terjadi, ketiga himpunan yang tersisa harus semuanya berasal
dari sisi kiri tabel atau semuanya dari sisi kanan
(lihat \nearbyexercise{exer:CancelPolSci}).
Ini menunjukkan adanya hubungan antara siklus mayoritas dan
penguraian yang kita gunakan\Dash paradoks
pemungutan suara hanya dapat terjadi ketika
kecenderungan-kecenderungan menuju preferensi siklik saling memperkuat.

Untuk pembuktiannya, andaikan kita telah membatalkan
urutan preferensi yang berlawanan dan tersisa satu himpunan
daftar preferensi untuk setiap baris.
Perhatikan daftar preferensi yang tersisa pada baris pertama.
Daftar itu dapat berasal dari sisi kiri atau kanan baris pertama
(atau tidak tersisa dari keduanya, jika daftar tersebut saling membatalkan dengan tepat).
Kita tuliskan
\begin{equation*}
  \votinggraphic{27}
\end{equation*}
dengan $a$ suatu bilangan bulat yang positif jika daftar yang tersisa berada di
sisi kiri, $a$ negatif jika daftar itu berada di sisi kanan, dan nol
jika pembatalannya sempurna.
Dengan cara serupa, kita mempunyai bilangan bulat $b$ dan $c$ untuk baris kedua dan ketiga;
masing-masing dapat bernilai positif, negatif, atau nol.

Hasil pemilihan kemudian ditentukan oleh jumlah berikut.
\begin{equation*}
  \votinggraphic{27}  %\votecycle{a}{a}{-a}{4} 
  +  
  \votinggraphic{28}  %\votecycle{-b}{b}{b}{4} 
  +  
  \votinggraphic{29}  %\votecycle{c}{-c}{c}{4}  
  =\hspace{.30in}
  \votinggraphic{30}  %\votecycle{a-b+c}{a+b-c}{-a+b+c}{4}
  \hspace{.30in}\hbox{}
\end{equation*}
Paradoks pemungutan suara terjadi ketika ketiga bilangan pada siklus total di sebelah kanan,
$-a+b+c$, $a-b+c$, dan $a+b-c$,
semuanya tak negatif atau semuanya tak positif.
Kita akan membuktikan bahwa hal ini hanya terjadi ketika
$a$, $b$, dan~$c$ semuanya tak negatif atau ketiganya tak positif.

Misalkan bilangan-bilangan pada siklus total tak negatif; kasus lainnya serupa.
\begin{equation*}
  \begin{linsys}{3}
    -a &+ &b &+ &c &\geq &0 \\
     a &- &b &+ &c &\geq &0 \\
     a &+ &b &- &c &\geq &0 
  \end{linsys}
\end{equation*}
Jumlahkan dua baris pertama untuk memperoleh $c\geq 0$.
Jumlahkan baris pertama dan ketiga untuk memperoleh $b\geq 0$.
Selanjutnya, penjumlahan baris kedua dan ketiga memberikan $a\geq 0$.
Jadi, jika siklus total tak negatif, maka pada setiap baris daftar
preferensi yang tersisa berasal dari sisi kiri tabel.
Dengan demikian, pembuktian selesai.

Hasil ini hanya menyatakan bahwa ketiga putaran yang searah merupakan
syarat perlu bagi siklus mayoritas.
Syarat itu tidak cukup; lihat
\nearbyexercise{exer:VoterCondNotSuff}.

Teori pemungutan suara dan topik-topik terkait merupakan bidang penelitian masa kini.
Ada banyak hasil yang menarik, terutama
hasil K~Arrow \cite{Arrow}, yang memperoleh Hadiah Nobel antara lain berkat
karya ini, yang menunjukkan bahwa tidak ada sistem pemungutan suara yang sepenuhnya adil
(untuk suatu definisi ``adil'' yang wajar).
Beberapa artikel pengantar yang baik ialah \cite{Gardner70},
\cite{Gardner74}, \cite{Gardner80}, dan \cite{NeimiRiker}.
\cite{Taylor} merupakan buku yang mudah dibaca.
Daftar panjang kasus dari sejarah politik Amerika modern dalam
\cite{GamingVote} menunjukkan bahwa paradoks-paradoks ini
secara rutin dimanipulasi dalam praktik.
(Di sisi lain, penelitian yang cukup baru menunjukkan bahwa
perhitungan cara memanipulasi pemilihan secara umum dapat menjadi tidak layak dilakukan,
tetapi hal ini berada di luar cakupan kita.)
Topik ini diadaptasi dari \cite{Zwicker}.
\emph{(Catatan Penulis:~Saya ingin berterima kasih kepada Profesor~Zwicker
atas diskusinya yang ramah dan mencerahkan.)}

\begin{exercises}
  \item \label{IrrIndVote}
    Berikut ini salah satu cara yang masuk akal bagi seorang pemilih untuk mempunyai preferensi siklik.
    Andaikan pemilih ini menilai setiap calon
    berdasarkan masing-masing dari tiga kriteria.
    \begin{exparts}
      \partsitem Buatlah tabel dengan baris berlabel `Demokrat',
         `Republik', dan `Ketiga', serta kolom berlabel
         `karakter', `pengalaman', dan `kebijakan'.
         Dalam setiap kolom, tempatkan satu
         calon sebagai yang paling disukai,
         calon lain di tengah, dan calon yang tersisa
         sebagai yang paling tidak disukai.
      \partsitem Dalam pemeringkatan ini, apakah calon Demokrat lebih disukai daripada calon Republik
         pada (setidaknya) dua dari tiga kriteria, atau sebaliknya?
         Apakah calon Republik lebih disukai daripada calon Ketiga?
      \partsitem Apakah tabel yang baru saja
         dibuat mempunyai urutan preferensi siklik?
         Jika tidak, buatlah tabel yang mempunyainya.
    \end{exparts}
    Jadi, seorang pemilih dapat mempunyai preferensi siklik di antara para calon.
    Akan tetapi, paradoks yang dijelaskan di atas ialah bahwa sekalipun setiap pemilih memiliki
    daftar preferensi lurus,
    preferensi siklik tetap dapat muncul pada kelompok secara keseluruhan.
    \begin{answer}
      Contoh ini menghasilkan urutan preferensi tak rasional bagi
      seorang pemilih.
      \smallskip
      \begin{center}
        \begin{tabular}{l|c|c|c|}
             \multicolumn{1}{c}{\ }
                  &\multicolumn{1}{c}{\textit{karakter}}
                  &\multicolumn{1}{c}{\textit{pengalaman}}
                  &\multicolumn{1}{c}{\textit{kebijakan}}   \\ 
          \cline{2-4}
         \textit{Demokrat}   &teratas &tengah  &terbawah \\
         \textit{Republik}   &tengah  &terbawah &teratas \\
         \textit{Ketiga}     &terbawah &teratas &tengah  \\
          \cline{2-4}
        \end{tabular}
      \end{center}
      \smallskip
      Calon Demokrat lebih baik daripada calon Republik dalam karakter dan
      pengalaman.
      Calon Republik mengungguli calon Ketiga dalam karakter dan kebijakan.
      Selanjutnya, calon Ketiga mengungguli calon Demokrat dalam pengalaman dan kebijakan.
    \end{answer}
  \item \label{VoteCalcTable} 
    Hitunglah nilai-nilai dalam tabel penguraian.
    \begin{answer}
      Pertama, bandingkan penguraian $D>R>T$ yang dibahas dalam Topik ini
      \begin{equation*}
        \votepreflist{-1}{1}{1}
          =\frac{1}{3}\cdot\votepreflist{1}{1}{1}
           +\frac{2}{3}\cdot\votepreflist{-1}{1}{0}
           +\frac{2}{3}\cdot\votepreflist{-1}{0}{1}
      \end{equation*}
      dengan penguraian pemilih $T>R>D$ yang berlawanan.
      \begin{equation*}
        \votepreflist{1}{-1}{-1}
          =d_1\cdot\votepreflist{1}{1}{1}
           +d_2\cdot\votepreflist{-1}{1}{0}
           +d_3\cdot\votepreflist{-1}{0}{1}
      \tag*{}\end{equation*}
      Jelas bahwa yang kedua merupakan negatif dari yang pertama, sehingga
      $d_1=-1/3$, $d_2=-2/3$, dan $d_3=-2/3$.
      Prinsip ini berlaku untuk setiap pasangan pemilih yang berlawanan, sehingga kita hanya perlu
      melakukan perhitungan bagi seorang pemilih dari baris kedua dan seorang pemilih dari
      baris ketiga.
      Untuk pemilih dengan putaran positif pada baris kedua,
      \begin{equation*}
        \begin{linsys}{3}
          c_1  &-  &c_2  &-  &c_3  &=  &1 \\
          c_1  &+  &c_2  &   &     &=  &1  \\
          c_1  &   &     &+  &c_3  &=  &-1  
        \end{linsys}
        \grstep[-\rho_1+\rho_3]{-\rho_1+\rho_2}
        \repeatedgrstep{(-1/2)\rho_2+\rho_3}
        \begin{linsys}{3}
          c_1  &-  &c_2  &-  &c_3      &=  &1 \\
               &   &2c_2 &+  &c_3      &=  &0  \\
               &   &     &   &(3/2)c_3 &=  &-2  
        \end{linsys}
      \tag*{}\end{equation*}
      memberikan $c_3=-4/3$, $c_2=2/3$, dan $c_1=1/3$.
      Untuk pemilih dengan putaran positif pada baris ketiga,
      \begin{equation*}
        \begin{linsys}{3}
          c_1  &-  &c_2  &-  &c_3  &=  &1 \\
          c_1  &+  &c_2  &   &     &=  &-1  \\
          c_1  &   &     &+  &c_3  &=  &1  
        \end{linsys}
        \grstep[-\rho_1+\rho_3]{-\rho_1+\rho_2}
        \repeatedgrstep{(-1/2)\rho_2+\rho_3}
        \begin{linsys}{3}
          c_1  &-  &c_2  &-  &c_3      &=  &1 \\
               &   &2c_2 &+  &c_3      &=  &-2  \\
               &   &     &   &(3/2)c_3 &=  &1  
        \end{linsys}
        \tag*{}\end{equation*}
        memberikan $c_3=2/3$, $c_2=-4/3$, dan $c_1=1/3$.
      \end{answer}
  \item \label{exer:CancelPolSci} 
    Lakukan pembatalan urutan preferensi yang berlawanan
    pada simulasi pemilihan kelas Ilmu Politik.
    Apakah semua preferensi yang tersisa berasal dari tiga entri kiri dalam
    tabel atau dari tiga entri kanan?
    \begin{answer}
      Simulasi pemilihan bersesuaian dengan tabel pada
      halaman~\pageref{table:Voting} sebagaimana ditunjukkan dalam tabel pertama,
      dan setelah pembatalan hasilnya adalah tabel kedua.
     \smallskip
      \begin{center}
         % \renewcommand{\arraystretch}{1.2}
         \begin{tabular}{|c|c|} \hline
           \textit{putaran positif}
           &\textit{putaran negatif}  \\ \hline
           \begin{tabular}{l}
             $D>R>T$        \\
             5 pemilih
           \end{tabular}
          &\begin{tabular}{l}
             $T>R>D$          \\
             2 pemilih
           \end{tabular}                  \\ \hline
          \begin{tabular}{l}
             $R>T>D$           \\
             8 pemilih
          \end{tabular}
         &\begin{tabular}{l}
             $D>T>R$             \\
             4 pemilih
           \end{tabular}                  \\ \hline
          \begin{tabular}{l}
            $T>D>R$             \\
            8 pemilih
          \end{tabular}
         &\begin{tabular}{l}
            $R>D>T$            \\
            2 pemilih
          \end{tabular}                  \\ \hline
        \end{tabular}
        \qquad
         \begin{tabular}{|c|c|} \hline
           \textit{putaran positif}
           &\textit{putaran negatif}  \\ \hline
           \begin{tabular}{l}
             $D>R>T$        \\
             3 pemilih
           \end{tabular}
          &\begin{tabular}{l}
             $T>R>D$          \\
             \textit{-dibatalkan-}
           \end{tabular}                  \\ \hline
          \begin{tabular}{l}
             $R>T>D$           \\
             4 pemilih
          \end{tabular}
         &\begin{tabular}{l}
             $D>T>R$             \\
             \textit{-dibatalkan-}
           \end{tabular}                  \\ \hline
          \begin{tabular}{l}
            $T>D>R$             \\
            6 pemilih
          \end{tabular}
         &\begin{tabular}{l}
            $R>D>T$            \\
            \textit{-dibatalkan-}
          \end{tabular}                  \\ \hline
        \end{tabular}
       \end{center}
       \smallskip
       Ketiganya berasal dari sisi tabel yang sama (sisi kiri),
       sebagaimana dinyatakan oleh hasil dalam Topik ini.
       Sekarang kita dapat menghitung hasilnya dengan menggunakan jumlah setelah pembatalan,
       \begin{equation*}
         % \psset{xunit=4pt,yunit=4pt,runit=4pt} 
         3\cdot \votinggraphic{31}
%          \pspicture[.4](-10.5,-4.2)(10,4) %\psgrid(-10,-10)(10,10)
%            \SpecialCoor
%            \rput(3;90){\small $D$}
%            \psarc{->}(0,0){3}{10}{70} %-rd->
%               \rput[lb](3.3;20){\scriptsize $1$ voter} 
%            \rput[B](3;210){\small $T$}
%            \psarc{->}(0,0){3}{220}{320} %-tr->
%               \rput[t](3.1;270){\scriptsize $1$ voter} 
%            \rput[B](3;330){\small $R$}
%            \psarc{->}(0,0){3}{110}{175} %-dt->
%               \rput[rb](3.3;160){\scriptsize $-1$ voter} 
%          \endpspicture
         +4\cdot\votinggraphic{32}
%          \pspicture[.4](-10,-4.2)(10,4) %\psgrid(-10,-10)(10,10)
%            \SpecialCoor
%            \rput(3;90){\small $D$}
%            \psarc{->}(0,0){3}{10}{70} %-rd->
%               \rput[lb](3.3;20){\scriptsize $-1$ voter} 
%            \rput[B](3;210){\small $T$}
%            \psarc{->}(0,0){3}{220}{320} %-tr->
%               \rput[t](3.1;270){\scriptsize $1$ voter} 
%            \rput[B](3;330){\small $R$}
%            \psarc{->}(0,0){3}{110}{175} %-dt->
%               \rput[rb](3.3;160){\scriptsize $1$ voter} 
%          \endpspicture
         +6\cdot\votinggraphic{33}
%          \pspicture[.4](-9,-4.2)(10,4) %\psgrid(-10,-10)(10,10)
%            \SpecialCoor
%            \rput(3;90){\small $D$}
%            \psarc{->}(0,0){3}{10}{70} %-rd->
%               \rput[lb](3.3;20){\scriptsize $1$ voter} 
%            \rput[B](3;210){\small $T$}
%            \psarc{->}(0,0){3}{220}{320} %-tr->
%               \rput[t](3.1;270){\scriptsize $-1$ voter} 
%            \rput[B](3;330){\small $R$}
%            \psarc{->}(0,0){3}{110}{175} %-dt->
%               \rput[rb](3.3;160){\scriptsize $1$ voter} 
%          \endpspicture
         = \votinggraphic{34}
%          \pspicture[.4](-9,-4.2)(10,4) %\psgrid(-10,-10)(10,10)
%            \SpecialCoor
%            \rput(3;90){\small $D$}
%            \psarc{->}(0,0){3}{10}{70} %-rd->
%               \rput[lb](3.3;20){\scriptsize $5$ voters} 
%            \rput[B](3;210){\small $T$}
%            \psarc{->}(0,0){3}{220}{320} %-tr->
%               \rput[t](3.1;270){\scriptsize $1$ voter} 
%            \rput[B](3;330){\small $R$}
%            \psarc{->}(0,0){3}{110}{175} %-dt->
%               \rput[rb](3.3;160){\scriptsize $7$ voters} 
%          \endpspicture
       \tag*{}\end{equation*}
       untuk memperoleh hasil yang sama.
    \end{answer}
  \item \label{exer:VoterCondNotSuff} 
    Syarat perlu bahwa
    paradoks pemungutan suara hanya dapat terjadi jika ketiga daftar preferensi yang tersisa
    setelah pembatalan mempunyai putaran yang sama bukanlah syarat cukup.
    \begin{exparts}
      % \partsitem Continuing the nonnegative total cycle case 
      %   considered in the proof, 
      %   use the two inequalities $0\leq a-b+c$ and $0\leq -a+b+c$
      %   to show that $\absval{a-b}\leq c$.
      % \partsitem Also show that $c\leq a+b$, and hence that
      %   $\absval{a-b}\leq c\leq a+b$. 
      \partsitem Berikan contoh pemungutan suara yang mempunyai siklus mayoritas,
        tetapi penambahan satu pemilih lagi dengan putaran yang sama menyebabkan siklus itu
        hilang.
      \partsitem Dapatkah hal sebaliknya terjadi; dapatkah penambahan seorang pemilih dengan
        putaran yang ``salah'' menyebabkan munculnya siklus?
      \partsitem Berikan suatu syarat yang perlu sekaligus cukup
        untuk memperoleh siklus mayoritas.
    \end{exparts}
    \begin{answer}
      \begin{exparts}
        % \partsitem We can rewrite the two as $-c\leq a-b$ and $-c\leq b-a$.
        %   Either $a-b$ or $b-a$ is nonpositive and so $-c\leq-\absval{a-b}$,
        %   as required.
        % \partsitem This is immediate from the supposition that $0\leq a+b-c$.
        \partsitem Contoh sederhana dimulai dengan pemilihan tanpa pemilih,
          yang mempunyai siklus mayoritas trivial, lalu menambahkan seorang pemilih mana pun.
          Contoh yang lebih menarik menggunakan simulasi
          pemilihan kelas Ilmu Politik dan menambahkan dua pemilih $T>D>R$
          (untuk memenuhi kriteria
          ``penambahan satu pemilih lagi'' dalam soal,
          kita dapat menambahkan mereka satu per satu).
          Para pemilih baru mempunyai putaran positif, yaitu
          putaran suara yang tersisa setelah pembatalan dalam simulasi
          pemilihan semula.
          Berikut adalah tabel pemilih yang dihasilkan dan, di sebelahnya,
          hasil pembatalan.
      \smallskip
      \begin{center}
         % \renewcommand{\arraystretch}{1.3}
         \begin{tabular}{|c|c|} \hline
           \textit{putaran positif}
           &\textit{putaran negatif}  \\ \hline
           \begin{tabular}{l}
             $D>R>T$        \\
             5 pemilih
           \end{tabular}
          &\begin{tabular}{l}
             $T>R>D$          \\
             2 pemilih
           \end{tabular}                  \\ \hline
          \begin{tabular}{l}
             $R>T>D$           \\
             8 pemilih
          \end{tabular}
         &\begin{tabular}{l}
             $D>T>R$             \\
             4 pemilih
           \end{tabular}                  \\ \hline
          \begin{tabular}{l}
            $T>D>R$             \\
            10 pemilih
          \end{tabular}
         &\begin{tabular}{l}
            $R>D>T$            \\
            2 pemilih
          \end{tabular}                  \\ \hline
        \end{tabular}
        \qquad
         \begin{tabular}{|c|c|} \hline
           \textit{putaran positif}
           &\textit{putaran negatif}  \\ \hline
           \begin{tabular}{l}
             $D>R>T$        \\
             3 pemilih
           \end{tabular}
          &\begin{tabular}{l}
             $T>R>D$          \\
             \textit{-dibatalkan-}
           \end{tabular}                  \\ \hline
          \begin{tabular}{l}
             $R>T>D$           \\
             4 pemilih
          \end{tabular}
         &\begin{tabular}{l}
             $D>T>R$             \\
             \textit{-dibatalkan-}
           \end{tabular}                  \\ \hline
          \begin{tabular}{l}
            $T>D>R$             \\
            8 pemilih
          \end{tabular}
         &\begin{tabular}{l}
            $R>D>T$            \\
            \textit{-dibatalkan-}
          \end{tabular}                  \\ \hline
        \end{tabular}
       \end{center}
       \smallskip
       Berikut adalah pemilihannya dengan menggunakan jumlah setelah pembatalan.
       \begin{equation*}
         % \psset{xunit=4pt,yunit=4pt,runit=4pt} 
         3\cdot\votinggraphic{35}
%          \pspicture[.4](-10.5,-4.2)(10,4) %\psgrid(-10,-10)(10,10)
%            \SpecialCoor
%            \rput(3;90){\small $D$}
%            \psarc{->}(0,0){3}{10}{70} %-rd->
%               \rput[lb](3.3;20){\scriptsize $1$ voter} 
%            \rput[B](3;210){\small $T$}
%            \psarc{->}(0,0){3}{220}{320} %-tr->
%               \rput[t](3.1;270){\scriptsize $1$ voter} 
%            \rput[B](3;330){\small $R$}
%            \psarc{->}(0,0){3}{110}{175} %-dt->
%               \rput[rb](3.3;160){\scriptsize $-1$ voter} 
%          \endpspicture
         +4\cdot\votinggraphic{36}
%          \pspicture[.4](-10,-4.2)(10,4) %\psgrid(-10,-10)(10,10)
%            \SpecialCoor
%            \rput(3;90){\small $D$}
%            \psarc{->}(0,0){3}{10}{70} %-rd->
%               \rput[lb](3.3;20){\scriptsize $-1$ voter} 
%            \rput[B](3;210){\small $T$}
%            \psarc{->}(0,0){3}{220}{320} %-tr->
%               \rput[t](3.1;270){\scriptsize $1$ voter} 
%            \rput[B](3;330){\small $R$}
%            \psarc{->}(0,0){3}{110}{175} %-dt->
%               \rput[rb](3.3;160){\scriptsize $1$ voter} 
%          \endpspicture
         +8\cdot\votinggraphic{37}
%          \pspicture[.4](-9,-4.2)(10,4) %\psgrid(-10,-10)(10,10)
%            \SpecialCoor
%            \rput(3;90){\small $D$}
%            \psarc{->}(0,0){3}{10}{70} %-rd->
%               \rput[lb](3.3;20){\scriptsize $1$ voter} 
%            \rput[B](3;210){\small $T$}
%            \psarc{->}(0,0){3}{220}{320} %-tr->
%               \rput[t](3.1;270){\scriptsize $-1$ voter} 
%            \rput[B](3;330){\small $R$}
%            \psarc{->}(0,0){3}{110}{175} %-dt->
%               \rput[rb](3.3;160){\scriptsize $1$ voter} 
%          \endpspicture
         =\votinggraphic{38}
%          \pspicture[.4](-9,-4.2)(10,4) %\psgrid(-10,-10)(10,10)
%            \SpecialCoor
%            \rput(3;90){\small $D$}
%            \psarc{->}(0,0){3}{10}{70} %-rd->
%               \rput[lb](3.3;20){\scriptsize $7$ voters} 
%            \rput[B](3;210){\small $T$}
%            \psarc{->}(0,0){3}{220}{320} %-tr->
%               \rput[t](3.1;270){\scriptsize $-1$ voter} 
%            \rput[B](3;330){\small $R$}
%            \psarc{->}(0,0){3}{110}{175} %-dt->
%               \rput[rb](3.3;160){\scriptsize $9$ voters} 
%          \endpspicture
       \tag*{}\end{equation*}
       Siklus mayoritas itu memang telah hilang.
      \partsitem Balikkan langkah pada soal sebelumnya.
         Artinya, tambahkan pada hasil soal sebelumnya seorang pemilih yang
         membatalkan pemilih yang membuat siklus tersebut hilang.
      \partsitem Salah satu syarat semacam itu ialah bahwa, setelah pembatalan,
         ketiganya tak negatif atau ketiganya tak positif,
         serta:~$\absval{c}\leq\absval{a+b}$, $\absval{b}\leq\absval{a+c}$,
         dan $\absval{a}\leq\absval{b+c}$.
         Syarat ini diperoleh dari diagram berikut.
         \begin{equation*}
           \votinggraphic{27}  %\votecycle{a}{a}{-a}{4} 
           +  
           \votinggraphic{28}  %\votecycle{-b}{b}{b}{4} 
           +  
           \votinggraphic{29}  %\votecycle{c}{-c}{c}{4}  
           =\hspace{.30in}
           \votinggraphic{30}  %\votecycle{a-b+c}{a+b-c}{-a+b+c}{4}
           \hspace{.30in}\hbox{}  
         \end{equation*}
%          \begin{center}
%            \votecycle{a}{a}{-a}{4} 
%            +  
%            \votecycle{-b}{b}{b}{4} 
%            +  
%            \votecycle{c}{-c}{c}{4}  
%            =\hspace{.30in}  
%            \votecycle{a-b+c}{a+b-c}{-a+b+c}{4}\hspace{.30in}\hbox{}  
%          \end{center}
      \end{exparts}
    \end{answer}
\item Pemilihan dengan satu pemilih tidak dapat mempunyai siklus mayoritas karena
    kita mensyaratkan bahwa daftar pemilih tersebut harus rasional.
    \begin{exparts}
      \partsitem Tunjukkan bahwa pemilihan dengan dua pemilih dapat mempunyai siklus mayoritas.
        (Kita menganggap preferensi kelompok sebagai siklus mayoritas jika ketiga
        total kelompok tak negatif atau jika ketiganya tak positif---artinya,
        kita mengizinkan adanya nilai nol dalam preferensi kelompok.)
      \partsitem Tunjukkan bahwa untuk setiap jumlah pemilih yang lebih besar daripada satu,
        terdapat pemilihan dengan sebanyak itu pemilih yang menghasilkan siklus
        mayoritas.
    \end{exparts}
    \begin{answer}
      \begin{exparts}
        \partsitem Pemilihan dengan dua pemilih dapat mempunyai siklus mayoritas melalui dua cara.
          Pertama, kedua pemilih dapat saling berlawanan,
          sehingga setelah pembatalan dihasilkan pemilihan trivial (dengan
          siklus mayoritas yang seluruh nilainya nol).
          Kedua, kedua pemilih dapat mempunyai putaran yang sama, tetapi berasal dari
          baris yang berbeda, seperti berikut.
          \begin{equation*}
            % \psset{xunit=4pt,yunit=4pt,runit=4pt} 
            1\cdot \votinggraphic{39}
%             \pspicture[.4](-10.5,-4.2)(10,4) %\psgrid(-10,-10)(10,10)
%               \SpecialCoor
%               \rput(3;90){\small $D$}
%               \psarc{->}(0,0){3}{10}{70} %-rd->
%                  \rput[lb](3.3;20){\scriptsize $1$ voter} 
%               \rput[B](3;210){\small $T$}
%               \psarc{->}(0,0){3}{220}{320} %-tr->
%                  \rput[t](3.1;270){\scriptsize $1$ voter} 
%               \rput[B](3;330){\small $R$}
%               \psarc{->}(0,0){3}{110}{175} %-dt->
%                  \rput[rb](3.3;160){\scriptsize $-1$ voter} 
%             \endpspicture
            +1\cdot\votinggraphic{40}
%             \pspicture[.4](-10,-4.2)(10,4) %\psgrid(-10,-10)(10,10)
%               \SpecialCoor
%               \rput(3;90){\small $D$}
%               \psarc{->}(0,0){3}{10}{70} %-rd->
%                  \rput[lb](3.3;20){\scriptsize $-1$ voter} 
%               \rput[B](3;210){\small $T$}
%               \psarc{->}(0,0){3}{220}{320} %-tr->
%                  \rput[t](3.1;270){\scriptsize $1$ voter} 
%               \rput[B](3;330){\small $R$}
%               \psarc{->}(0,0){3}{110}{175} %-dt->
%                  \rput[rb](3.3;160){\scriptsize $1$ voter} 
%             \endpspicture
            +0\cdot\votinggraphic{41}
%             \pspicture[.4](-9,-4.2)(10,4) %\psgrid(-10,-10)(10,10)
%               \SpecialCoor
%               \rput(3;90){\small $D$}
%               \psarc{->}(0,0){3}{10}{70} %-rd->
%                  \rput[lb](3.3;20){\scriptsize $1$ voter} 
%               \rput[B](3;210){\small $T$}
%               \psarc{->}(0,0){3}{220}{320} %-tr->
%                  \rput[t](3.1;270){\scriptsize $-1$ voter} 
%               \rput[B](3;330){\small $R$}
%               \psarc{->}(0,0){3}{110}{175} %-dt->
%                  \rput[rb](3.3;160){\scriptsize $1$ voter} 
%             \endpspicture
            =\votinggraphic{42}
%             \pspicture[.4](-9,-4.2)(10,4) %\psgrid(-10,-10)(10,10)
%               \SpecialCoor
%               \rput(3;90){\small $D$}
%               \psarc{->}(0,0){3}{10}{70} %-rd->
%                  \rput[lb](3.3;20){\scriptsize $0$ voters} 
%               \rput[B](3;210){\small $T$}
%               \psarc{->}(0,0){3}{220}{320} %-tr->
%                  \rput[t](3.1;270){\scriptsize $2$ voters} 
%               \rput[B](3;330){\small $R$}
%               \psarc{->}(0,0){3}{110}{175} %-dt->
%                  \rput[rb](3.3;160){\scriptsize $0$ voters} 
%             \endpspicture
          \tag*{}\end{equation*}
        \partsitem Ada dua kasus.
          Pemilih yang berjumlah genap dapat dibagi dua sama banyak menjadi pasangan berlawanan;
          misalnya, separuh pemilih memilih $D>R>T$ dan separuh lainnya $T>R>D$.
          Pembatalan kemudian menghasilkan pemilihan trivial.
          Jika jumlah pemilih lebih besar daripada satu dan ganjil (berbentuk
          $2k+1$ dengan $k>0$), maka dengan menggunakan diagram siklus dari pembuktian,
          \begin{equation*}
            \votinggraphic{27}  %\votecycle{a}{a}{-a}{4} 
            +  
            \votinggraphic{28}  %\votecycle{-b}{b}{b}{4} 
            +  
            \votinggraphic{29}  %\votecycle{c}{-c}{c}{4}  
            =\hspace{.30in}
            \votinggraphic{30}  %\votecycle{a-b+c}{a+b-c}{-a+b+c}{4}
            \hspace{.30in}\hbox{}  
          \end{equation*}
%          \begin{center}
%            \votecycle{a}{a}{-a}{4} 
%            +  
%            \votecycle{-b}{b}{b}{4} 
%            +  
%            \votecycle{c}{-c}{c}{4}  
%            =\hspace{.30in}  
%            \votecycle{a-b+c}{a+b-c}{-a+b+c}{4}\hspace{.30in}\hbox{}  
%          \end{center}
          kita dapat mengambil $a=k$, $b=k$, dan $c=1$.
          Karena $k>0$, hasilnya merupakan siklus mayoritas.
       \end{exparts}
    \end{answer}
\item \label{exer:PerpIsSubSp} 
    Misalkan $U$ suatu subruang dari $\Re^3$.
    Buktikan bahwa himpunan
    $U^\perp=\set{\vec{v}\suchthat 
              \text{$\vec{v}\dotprod\vec{u}=0$ untuk setiap $\vec{u}\in U$}}$
    yang terdiri atas vektor-vektor
    yang tegak lurus terhadap setiap vektor dalam $U$ juga merupakan subruang dari $\Re^3$.
    Apakah pernyataan ini tetap berlaku jika $U$ bukan subruang?
    \begin{answer}
      Himpunan ini tak kosong karena memuat vektor nol.
      Untuk melihat bahwa himpunan ini tertutup terhadap kombinasi linear dua
      anggotanya, andaikan $\vec{v}_1$ dan $\vec{v}_2$ berada dalam $U^\perp$
      dan perhatikan $c_1\vec{v}_1+c_2\vec{v}_2$.
      Untuk setiap $\vec{u}\in U$,
      \begin{equation*}
        (c_1\vec{v}_1+c_2\vec{v}_2)\dotprod\vec{u}
        =c_1(\vec{v}_1\dotprod\vec{u})+c_2(\vec{v}_2\dotprod\vec{u})
        =c_1\cdot 0+c_2\cdot 0=0         
      \tag*{}\end{equation*}
      sehingga $c_1\vec{v}_1+c_2\vec{v}_2\in U^\perp$.

      Mengenai apakah pernyataan tersebut tetap berlaku jika $U$ merupakan himpunan bagian,
      tetapi bukan subruang, jawabannya adalah ya.
    \end{answer}
\end{exercises}
\index{paradoks pemungutan suara|)}
\endinput
